\documentclass[12pt, twoside]{article}
\input{../preamble}

\fancyhead[LE]{\thepage}
\fancyhead[RO]{Markscheme}
\fancyhead[LO]{La Scuola d'Italia / Dr.\ Huson / IB Math \\* 18 April 2026}

\begin{document}
\subsubsection*{5.14 Classwork: Logarithms and Exponentials --- Markscheme}

\begin{enumerate}[itemsep=0.3cm]

%--- Problem 1: Product rule with exponential [5 marks] ---
%--- Source: AA SL 2022 May P1 TZ1 Q05 ---
\item[1.] [Maximum mark: 5]

Consider the curve with equation $y = (2x - 1)e^{kx}$,
where $x \in \mathbb{R}$ and $k \in \mathbb{R}$.
The tangent to the curve at the point where $x = 1$
is parallel to the line $y = 5xe^k$.
Find the value of $k$.

\textbf{Solution:}

Product rule: \hfill \textbf{(M1)}

$\dfrac{dy}{dx} = 2 \cdot e^{kx} + (2x - 1) \cdot ke^{kx}
= e^{kx}(2kx - k + 2)$ \hfill \textbf{(A1)}

At $x = 1$:

$\dfrac{dy}{dx} = e^k(k + 2)$ \hfill \textbf{(A1)}

The line $y = 5xe^k$ has slope $5e^k$.

Setting equal: \hfill \textbf{(M1)}

$e^k(k + 2) = 5e^k$

$k + 2 = 5$

$k = 3$ \hfill \textbf{(A1)}

\bigskip
\hrule
\bigskip

%--- Problem 2: Logarithmic function [9 marks] ---
%--- Source: AA SL 2021 May P1 TZ2 Q05 ---
\item[2.] [Maximum mark: 9]

Consider the function $f$ defined by
$f(x) = \ln(x^2 - 16)$ for $x > 4$.

\medskip

\textbf{(a)} Find the exact value of $a$. \hfill [3]

\textbf{Solution:}

$f(a) = 0 \implies \ln(a^2 - 16) = 0$ \hfill \textbf{(M1)}

$a^2 - 16 = 1$ \hfill \textbf{(A1)}

$a^2 = 17 \implies a = \sqrt{17}$ \hfill \textbf{(A1)}

\medskip

\textbf{(b)} Given that the gradient of $L$ is $\dfrac{1}{3}$,
find the $x$-coordinate of $B$. \hfill [6]

\textbf{Solution:}

$f'(x) = \dfrac{2x}{x^2 - 16}$ \hfill \textbf{(M1)(A1)}

Setting $f'(x) = \dfrac{1}{3}$: \hfill \textbf{(M1)}

$\dfrac{2x}{x^2 - 16} = \dfrac{1}{3}$

$6x = x^2 - 16$ \hfill \textbf{(A1)}

$x^2 - 6x - 16 = 0$

$(x - 8)(x + 2) = 0$ \hfill \textbf{(A1)}

$x = 8$ \quad (since $x > 4$) \hfill \textbf{(A1)}

\newpage

%--- Problem 3: Exponential function and inverse [14 marks] ---
%--- Source: Math SL 2019 Nov P1 Q10 ---
\item[3.] [Maximum mark: 14]

Let $g(x) = p^x + q$, for $x, \, p, \, q \in \mathbb{R}$, $p > 1$.
The point $A(0, a)$ lies on the graph of $g$.
Let $f(x) = g^{-1}(x)$.

\medskip

\textbf{(a)} Write down the coordinates of $B$. \hfill [2]

\textbf{Solution:}

$g(0) = p^0 + q = 1 + q$, so $a = q + 1$ \hfill \textbf{(A1)}

$B$ is the reflection of $A(0, a)$ in $y = x$,
so $B = (a, 0)$ \hfill \textbf{(A1)}

\medskip

\textbf{(b)} Given that $f'(a) = \dfrac{1}{\ln p}$, find the
equation of $L_1$ in terms of $x$, $p$ and $q$. \hfill [5]

\textbf{Solution:}

$L_1$ passes through $B(a, 0)$ with gradient
$\dfrac{1}{\ln p}$ \hfill \textbf{(M1)}

$y - 0 = \dfrac{1}{\ln p}(x - a)$ \hfill \textbf{(A1)}

Substituting $a = q + 1$: \hfill \textbf{(M1)(A1)}

$y = \dfrac{1}{\ln p}(x - q - 1)$ \hfill \textbf{(A1)}

\medskip

\textbf{(c)} Find the equation of $L_1$ in terms of $x$. \hfill [7]

\textbf{Solution:}

\emph{Finding $p$:}

The gradient of the normal to $g$ at $A$ is
$\dfrac{1}{\ln \frac{1}{3}} = \dfrac{1}{-\ln 3}
= -\dfrac{1}{\ln 3}$ \hfill \textbf{(M1)}

Since (gradient of tangent) $\times$ (gradient of normal) $= -1$:

$(\ln p) \cdot \left(-\dfrac{1}{\ln 3}\right) = -1$ \hfill \textbf{(A1)}

$\ln p = \ln 3 \implies p = 3$ \hfill \textbf{(A1)}

\emph{Finding $q$:}

$L_2$ passes through $(-2, -2)$:

$-2 = (\ln 3)(-2) + q + 1$ \hfill \textbf{(A1)}

$q = -3 + 2\ln 3$ \hfill \textbf{(A1)}

\emph{Finding $L_1$:}

$a = q + 1 = -2 + 2\ln 3$ \hfill \textbf{(A1)}

$L_1: \; y = \dfrac{1}{\ln 3}(x - 2\ln 3 + 2)$ \hfill \textbf{(A1)}

\end{enumerate}
\end{document}
